\begin{explanationbox}
	\textbf{\textcolor{CExplanation}{一句话直觉}}

	二次三项式配方后变为$(x+\frac{b}{2a})^2$乘以系数再加常数，从而直接看出对称轴和最值。

	\medskip
	\textbf{\textcolor{CExplanation}{核心拆解}}
	\begin{description}[style=nextline, leftmargin=2.8em, labelsep=0.5em, itemsep=0.45em]
		\item[\textbf{要点一（配方结构）：}] 将二次项和一次项组合成完全平方 $a(x+\frac{b}{2a})^2$。
		\item[\textbf{要点二（常数分离）：}] 余下的常数 $\frac{4ac-b^2}{4a}$ 是配方后的修正项。
	\end{description}

	\medskip
	\textbf{\textcolor{CExplanation}{几何本质（顶点形式）}}

	配方后的表达式 $a(x+\frac{b}{2a})^2 + \frac{4ac-b^2}{4a}$ 就是顶点式，顶点为 $(-\frac{b}{2a},\frac{4ac-b^2}{4a})$。

	\medskip
	\textbf{\textcolor{CExplanation}{代数意义（降次工具）}}

	通过配方将二次式 $ax^2+bx+c$ 写成平方形式，从而可以开平方求解方程，或进行不等式放缩。

	\medskip
	\textbf{\textcolor{CExplanation}{考点价值（命题热点）}}
	\begin{itemize}[leftmargin=*, itemsep=0.5em, topsep=0.4em]
		\item \textbf{考法一（最值问题）：} 已知二次函数解析式，求最大值或最小值。
		\item \textbf{考法二（解二次方程）：} 利用配方后开平方得到求根公式。
		\item \textbf{考法三（不等式证明）：} 配方后利用平方非负性证明二次不等式。
	\end{itemize}

	\medskip
	\textbf{\textcolor{CExplanation}{顿悟点}}

	配方后自动给出对称轴 $x=-\frac{b}{2a}$ 和最值 $\frac{4ac-b^2}{4a}$，无需额外计算。

	\medskip
	\textbf{\textcolor{CExplanation}{使用场景}}
	\begin{itemize}[leftmargin=*, itemsep=0.5em, topsep=0.4em]
		\item \textbf{场景一（求顶点坐标）：} 需要二次函数图像性质时直接配方。
		\item \textbf{场景二（解二次方程）：} 当二次方程不易因式分解时，优先配方。
		\item \textbf{场景三（证明不等式）：} 将二次式配方后利用平方非负性。
	\end{itemize}
\end{explanationbox}
